\section{两大模型的对比与联系}

\subsection{VAE vs. GAN：架构对比}

\begin{figure}[H]
\centering
\includegraphics[width=0.85\textwidth]{slides-images/slide-46.jpg}
\caption{两大隐变量模型家族：VAE 和 GAN\protect\footnotemark}
\end{figure}
\footnotetext{Slides 第46页。}

\begin{table}[H]
\centering
\caption{VAE 与 GAN 的系统对比}
\begin{tabular}{lll}
\toprule
\textbf{维度} & \textbf{VAE} & \textbf{GAN} \\
\midrule
建模方式 & 显式建模密度 $p(x)$ & 隐式学习采样过程 \\
网络结构 & 编码器 + 解码器 & 生成器 + 判别器 \\
训练方式 & 最大化 ELBO & Minimax 对抗博弈 \\
隐空间 & 编码器产生，结构化 & 噪声空间，无编码器 \\
生成质量 & 通常较模糊 & 通常更清晰锐利 \\
训练稳定性 & 相对稳定 & 不稳定，易模式崩塌 \\
可解释性 & 隐空间可解释 & 隐空间解释性较弱 \\
\bottomrule
\end{tabular}
\end{table}

\begin{warningbox}{VAE 生成模糊的原因}
VAE 使用像素级重建损失（如 MSE），这导致模型倾向于输出多种可能结果的"平均"，产生模糊的输出。GAN 则通过判别器提供的对抗信号来优化生成质量，不受这个限制，因此通常能生成更清晰的图像。但 GAN 缺少 VAE 那样的显式编码器，难以进行推断和隐空间操控。
\end{warningbox}

\subsection{本章小结}

VAE 和 GAN 是两种互补的生成模型范式。VAE 提供了结构化的隐空间和稳定的训练，但生成质量有限；GAN 提供了卓越的生成质量，但训练不稳定且缺乏显式的隐空间。在实际应用中，可以根据需求选择合适的模型，或使用结合两者优点的混合架构。


